\documentclass[17pt]{article}

\usepackage{cmap}	% поиск по документу
\usepackage{hyperref}	% поддержка гиперссылок
\usepackage{geometry}	% настройка геометрии документа
\usepackage{multicol}   % настройка шаблона расположения текста
\usepackage{extsizes}	% расширенная настройка размеров шрифтов
\usepackage{indentfirst}    % отступ для первого абзаца
\usepackage{tabularx}	% расширенные таблицы
\usepackage{makecell}   % настройка клеток таблиц
\usepackage{caption}    % настройка подписей к вставкам
\usepackage[table]{xcolor}  % покраска таблиц
\usepackage{graphicx}	% вставка изображений
\usepackage{wrapfig, booktabs}	% оформление таблиц
\usepackage{graphicx, caption, subcaption}
\usepackage{minted}    % вставка кода

\usepackage{color}	% настройка цвета текста - \textcolor
\usepackage{soul}	% выделение текста - \hl
\usepackage[T2A]{fontenc}	% поддержка кириллических символов
\usepackage[utf8]{inputenc} % поддержка кириллических символов
\usepackage[english,russian]{babel}	% поддержка кириллических символов
\usepackage{csquotes, biblatex}   % вывод библиографии

\usepackage{mathcmd}	% математические команды
\usepackage{amsfonts}	% математические символы 1
\usepackage{amssymb}	% математические символы 2
\usepackage{amsmath}	% математические символы 3
\usepackage{amsthm}	% работа с теоремами

% Настройка окружения теорем 1 (RU)
\newtheorem{theorem}{Теорема}[section]
\newtheorem{corollary}{Следствие}[theorem]
\newtheorem{lemma}[theorem]{Лемма}
\newtheorem{example}{Пример}[section]
\newtheorem*{remark}{Утверждение}
\addto\captionsrussian{\renewcommand\proofname{Док-во:}}

% Настройка окружения теорем 2 (EN)
%\newtheorem{theorem}{TH}[section]
%\newtheorem{corollary}{CO}[theorem]
%\newtheorem{lemma}[theorem]{LM}
%\newtheorem{example}{EX}[section]
%\addto\captionsrussian{\renewcommand\proofname{PF:}}

\renewcommand{\labelitemii}{$\circ$}
\renewcommand{\labelitemiii}{-}

\addbibresource{biblio.bib}

% Настройка оформления вставок кода
\setminted{xleftmargin=\parindent, fontsize=\footnotesize, linenos, tabsize=2, breaklines}
\usemintedstyle{emacs} 

% Настройка отступов
\geometry{
    a4paper,
    % total={200mm,287mm},
    left=10mm,
    right=10mm,
    top=10mm,
    bottom=20mm
}

% Заголовок
\title{
    \textbf{Методы моделирования и анализа стохастических систем} \\
    \large Отчет по лабораторным работам
}
\author{Костарев Иван (МЕНМ-150201)}
\date{весна 2025/2026}

\begin{document}

\maketitle
\tableofcontents

\graphicspath{ {./lab1/images/part1/} }

\newpage
\section{Лабораторная работа № 1, часть 1: Генерация псевдослучайных чисел}
\subsection{Инициализация ГПСЧ для X $\sim$ U[0, 1]}

Инициализируем генератор \texttt{rng}, используя конструктор \texttt{default\_rng} из библиотеки \texttt{numpy}:
\begin{minted}{python}
rng = np.random.default_rng()
\end{minted}



\subsection{Нахождение зависимости между мат ожиданием, дисперсией и размером выборки}

Для исследования сходимости выборочных характеристик генерируются массивы различной длины и вычисляются их математические ожидания и дисперсии.

Построим график (Рис. \ref{fig:mean_vs_N}) зависимости математического ожидания \texttt{mean} от размера выборки $N$.
\begin{figure}[htbp]
    \centering
    \includegraphics[width=0.6\textwidth]{mean_vs_N.png}
    \caption{Зависимость выборочного математического ожидания от размера выборки $N$.}
    \label{fig:mean_vs_N}
\end{figure}

Построим график (Рис. \ref{fig:var_vs_N}) зависимости дисперсии \texttt{var} от размера выборки $N$.
\begin{figure}[htbp]
    \centering
    \includegraphics[width=0.6\textwidth]{var_vs_N.png}
    \caption{Зависимость выборочной дисперсии от размера выборки $N$.}
    \label{fig:var_vs_N}
\end{figure}



\newpage
\subsection{Построение гистограммы и вычисление погрешности (расхождение сгенерированных чисел с теоретическим значением)}

Генерируется выборка большого объема ($N = 10^8$) и строится ее нормированная гистограмма:
\begin{minted}{python}
sample_size = 10 ** 8
sample = rng.uniform(0., 1., size=sample_size)
num_bins = 32
sample_hist, bin_edges = np.histogram(sample, bins=num_bins, density=True)
\end{minted}

Построим приближенную гистограмму (Рис. \ref{fig:histogram}) равномерного распределения на основе сгенерированных данных (\texttt{sample}).
\begin{figure}[htbp]
    \centering
    \includegraphics[width=0.6\textwidth]{histogram.png}
    \caption{Приближенная гистограмма равномерного распределения $U[0, 1]$.}
    \label{fig:histogram}
\end{figure}

Определим функцию для вычисления погрешности приближения гистограммы к теоретической плотности ($p_{\text{true}} = 1$):
\begin{minted}{python}
def distribution_error(p: np.ndarray, p_true: float = 1.):
    return np.sum((p - p_true) ** 2)
\end{minted}

Построим график (Рис. \ref{fig:error_vs_N}) зависимости погрешности \texttt{distribution\_error} от размера выборки $N$.
\begin{figure}[htbp]
    \centering
    \includegraphics[width=0.6\textwidth]{error_vs_N.png}
    \caption{Зависимость погрешности распределения от размера выборки $N$.}
    \label{fig:error_vs_N}
\end{figure}



\newpage
\subsection{Вычисление коэффициента корреляции между двумя выборками из одного распределения, но построенными разными способами}

Для оценки линейной зависимости реализуем вычисление коэффициента корреляции Пирсона:
\begin{minted}{python}
def pearson_r(x: np.ndarray, y: np.ndarray):
    return np.sum((x - x.mean()) * (y - y.mean())) / np.sqrt(np.sum((x - x.mean()) ** 2) * np.sum((y - y.mean()) ** 2))
\end{minted}

\vspace{0.5em}
\textbf{(а)} генерация двух независимых выборок.

Построим график (Рис. \ref{fig:corr_a}) зависимости коэффициента корреляции \texttt{pearson\_r} от размера выборки $N$ (вариант а).
\begin{figure}[htbp]
    \centering
    \includegraphics[width=0.6\textwidth]{corr_a.png}
    \caption{Зависимость коэффициента корреляции от размера выборки $N$ (вариант а).}
    \label{fig:corr_a}
\end{figure}

\vspace{0.5em}
\textbf{(б)} генерация единой последовательности и разделение на четные/нечетные элементы.

Построим график (Рис. \ref{fig:corr_b}) зависимости коэффициента корреляции \texttt{pearson\_r} от размера выборки $N$ (вариант б).
\begin{figure}[htbp]
    \centering
    \includegraphics[width=0.6\textwidth]{corr_b.png}
    \caption{Зависимость коэффициента корреляции от размера выборки $N$ (вариант б).}
    \label{fig:corr_b}
\end{figure}

\graphicspath{ {./lab1/images/part2/} }

\newpage
\section{Лабораторная работа № 1, часть 2: Равномерное распределение}
\subsection{Метод обратной функции}

Задана плотность

$$
\begin{cases}
    f(x) = \frac{e^x}{e - 1}, & x \in (0, 1) \\
    0, & \text{иначе}
\end{cases}
$$

Из плотности $f(x)$ найдем распределение $F(x)$:

$$
F(x) = \int_{-\infty}^{x} f(t)dt = \int_{-\infty}^{0} 0 dt + \int_{0}^{x} \frac{e^t}{e - 1}dt = \frac{1}{e - 1}(e^x - e^{0}) = \frac{e^x - 1}{e - 1}
$$
при $x \in (0, 1)$.

Найдем обратную функцию $G(y) = F^{-1}(y)$:

$$
G(y) = \ln{1 + y(e - 1)}
$$
при $y \in (0, 1)$.

\textbf{(а)} Смоделируем $N = 10^8$ значений величины из распределения $F(x)$. Для этого используется функция $G(y)$, примененная к равномерно распределенной выборке.

\begin{minted}[fontsize=\footnotesize]{python}
def G(y):
    return np.log(1 + y * (np.e - 1))

sample_size = 10 ** 8
sample = rng.uniform(0., 1., size=sample_size)
sample = G(sample)
\end{minted}

На Рис. \ref{fig:hist_inv} представлена гистограмма смоделированной выборки.

\begin{figure}[h!]
    \centering
    \includegraphics[width=0.6\textwidth]{hist_inv.png}
    \caption{Гистограмма выборки $X \sim F(x)$ на интервале $[0, 1]$}
    \label{fig:hist_inv}
\end{figure}

\textbf{(б)} Найдем точные значения мат ожидания $M_x$ и дисперсии $D_x$.

$$
M_x = \int_{0}^{1} x f(x)dx = \int_{0}^{1} x \frac{e^x}{e - 1}dx = \frac{1}{e - 1}
$$

$$
D_x = \int_{0}^{1} x^2 f(x)dx - M_x^2 = \int_{0}^{1} x^2 \frac{e^x}{e - 1}dx - M_x^2 = \frac{e - 2}{e - 1} - (\frac{1}{e - 1})^2 = \frac{e^2 - 3e + 1}{(e - 1)^2}.
$$

\textbf{(в)} Вычислим оценки мат ожидания и дисперсии для различных размеров выборки $N$ и сравним их с точными значениями (Рис. \ref{fig:mean_inv}, \ref{fig:var_inv}).

\begin{figure}[h!]
    \centering
    \includegraphics[width=0.6\textwidth]{mean_comp_inv.png}
    \caption{Сравнение выборочного среднего с точным значением}
    \label{fig:mean_inv}
\end{figure}

\begin{figure}[h!]
    \centering
    \includegraphics[width=0.6\textwidth]{var_comp_inv.png}
    \caption{Сравнение выборочной дисперсии с точным значением}
    \label{fig:var_inv}
\end{figure}

\textbf{Итог}: при размере выборки $N \geq 10^5$ оценки мат ожидания и дисперсии слабо отличаются от точных значений.



\newpage
\subsection{Метод Монте-Карло}

Задана функция $f(x) = \sqrt{x} + 1$, $x \in (1, 4)$.

Найдем точное значение интеграла:

$$
I = \int_{1}^{4} f(x)dx = \int_{1}^{4} (\sqrt{x} + 1)dx = \frac{23}{3} \approx 7.67
$$

Найдем приближенное значение интеграла методом Монте-Карло и сравним его с точным (Рис. \ref{fig:mc_comp}).

\begin{figure}[h!]
    \centering
    \includegraphics[width=0.6\textwidth]{mc_comp.png}
    \caption{Сравнение оценки интеграла с точным значением}
    \label{fig:mc_comp}
\end{figure}

Визуализируем метод Монте-Карло для функции $f(x)$ на примере $N = 1000$ случайных точек (Рис. \ref{fig:mc_vis}).

\begin{figure}[h!]
    \centering
    \includegraphics[width=0.6\textwidth]{mc_vis.png}
    \caption{Визуализация метода Монте-Карло для $N=1000$}
    \label{fig:mc_vis}
\end{figure}

\textbf{Итог}: при количестве точек $N \geq 10^7$ оценка интеграла слабо отличается от точного значения.

\graphicspath{ {./lab2/images/} }

\newpage
\section{Лабораторная работа № 2: Нормальное распределение}
\subsection{Метод 12}

Программная реализация метода 12:

\begin{minted}{python}
def random_normal_method_of_12(size: int = 1):
    n = 12
    uniform = rng.uniform(0., 1., size=(n, size)) - 0.5
    return np.sum(uniform, axis=0) * np.sqrt(12) / np.sqrt(n)    
\end{minted}

Смоделируем методом 12 нормально распределенную выборку $X \sim \mathcal{N}(0, 1)$ размера $N = 10^7$, и построим гистограмму плотности распределения (Рис. \ref{fig:method12_hist}).

\begin{figure}[htbp]
    \centering
    \includegraphics[width=0.6\textwidth]{method12_hist.png}
    \caption{Эмпирическая плотность распределения величины, полученной методом 12.}
    \label{fig:method12_hist}
\end{figure}

Программная реализация функции ошибки:

\begin{minted}{python}
def distribution_error(
    p: np.ndarray,
    p_true: np.ndarray
):
    return np.sum((p - p_true) ** 2)  
\end{minted}

Построим график зависимости погрешности \texttt{distribution\_error} от размера выборки $N$ (Рис. \ref{fig:method12_error}).

\begin{figure}[htbp]
    \centering
    \includegraphics[width=0.6\textwidth]{method12_error.png}
    \caption{Зависимость погрешности гистограммы от размера выборки $N$.}
    \label{fig:method12_error}
\end{figure}

Вычислим оценки мат ожидания и дисперсии для различных размеров выборки $N$, и сравним их с точными значениями для распределения $\mathcal{N}(0, 1)$ (Рис. \ref{fig:method12_means_plot}, \ref{fig:method12_vars_plot}).

\begin{figure}[htbp]
    \centering
    \includegraphics[width=0.6\textwidth]{method12_means_plot.png}
    \caption{Сравнение выборочного среднего с точным значением мат ожидания.}
    \label{fig:method12_means_plot}
\end{figure}

\begin{figure}[htbp]
    \centering
    \includegraphics[width=0.6\textwidth]{method12_vars_plot.png}
    \caption{Сравнение выборочной дисперсии с точным значением дисперсии.}
    \label{fig:method12_vars_plot}
\end{figure}

\textbf{Итог}: метод 12 для генерации нормально распределенной выборки требует довольно много памяти в текущей реализации на Python, поэтому эмпирические параметры распределения были вычислены для выборок размеров $N \leq 10^7$. Даже на таких размерах выборок графики демонстрируют, что с увеличением размера выборки мат ожидание и дисперсия стремятся к точным значениям. Ошибка плотности по сравнению с истинным нормальным распределением стремится к 0.



\newpage
\subsection{Преобразование Бокса -- Мюллера}

Программная реализация преобразования Бокса -- Мюллера:

\begin{minted}{python}
def random_normal_box_muller(size: int = 1):
    a, b = rng.uniform(0., 1., size=(2, size))
    xi_1 = np.sqrt(-2 * np.log(a)) * np.cos(2 * np.pi * b)
    xi_2 = np.sqrt(-2 * np.log(a)) * np.sin(2 * np.pi * b)
    return xi_1
\end{minted}

Смоделируем с помощью преобразования Бокса -- Мюллера нормально распределенную выборку $X \sim \mathcal{N}(0, 1)$ размера $N = 10^7$, и построим гистограмму плотности распределения (Рис. \ref{fig:box_muller_hist}).

\begin{figure}[htbp]
    \centering
    \includegraphics[width=0.6\textwidth]{box_muller_hist.png}
    \caption{Эмпирическая плотность распределения величины, полученной преобразованием Бокса -- Мюллера.}
    \label{fig:box_muller_hist}
\end{figure}

Построим график зависимости погрешности \texttt{distribution\_error} от размера выборки $N$ (Рис. \ref{fig:box_muller_error}).

\begin{figure}[htbp]
    \centering
    \includegraphics[width=0.6\textwidth]{box_muller_error.png}
    \caption{Зависимость погрешности гистограммы от размера выборки $N$.}
    \label{fig:box_muller_error}
\end{figure}

Вычислим оценки мат ожидания и дисперсии для различных размеров выборки $N$, и сравним их с точными значениями для распределения $\mathcal{N}(0, 1)$.

\begin{figure}[htbp]
    \centering
    \includegraphics[width=0.6\textwidth]{box_muller_means_plot.png}
    \caption{Сравнение выборочного среднего с точным значением мат ожидания.}
    \label{fig:box_muller_means_plot}
\end{figure}

\begin{figure}[htbp]
    \centering
    \includegraphics[width=0.6\textwidth]{box_muller_vars_plot.png}
    \caption{Сравнение выборочной дисперсии с точным значением дисперсии.}
    \label{fig:box_muller_vars_plot}
\end{figure}

\textbf{Итог}: преобразование Бокса -- Мюллера для генерации нормально распределенной выборки требует довольно много памяти в текущей реализации на Python, поэтому эмпирические параметры распределения были вычислены для выборок размеров $N \leq 10^7$. Даже на таких размерах выборок графики демонстрируют, что с увеличением размера выборки мат ожидание и дисперсия стремятся к точным значениям. Ошибка плотности по сравнению с истинным нормальным распределением стремится к 0.

\graphicspath{ {./lab3/images/part1/} }

\newpage
\section{Лабораторная работа № 3, часть 1: Анализ детерминированной задачи}
\subsection{Нейронная модель ФитцХью-Нагумо (FHN)}

$$
\begin{cases}
    \dot{x} = \frac{1}{\delta} (x - \frac{x^3}{3} - y) \\
    \dot{y} = x + a
\end{cases}
$$

1. Найдем равновесия системы.

$$
\begin{cases}
    \frac{1}{\delta} (x - \frac{x^3}{3} - y) = 0 \\
    x + a = 0
\end{cases} \implies
\begin{cases}
    \overline{x} = -a \\
    \overline{y} = \frac{a^3}{3} - a
\end{cases}
$$

2. Исследуем равновесия на устойчивость.

$$
J(x, y) = \begin{pmatrix}
    \frac{1}{\delta} (1 - x^2) & \frac{1}{\delta} \\
    1 & 0
\end{pmatrix}
$$

$$
|J(x, y) - \lambda E| = 0
$$

$$
\lambda_{1,2}(x, y) = \frac{-\delta x^2 + \delta \pm \delta \sqrt{x^4 - 2 x^2 - 4 \delta + 1}}{2 \delta^2}
$$

Далее с помощью численной процедуры исследуем поведение собственных чисел якобиана в равновесных точках при различных значениях параметра $a$.

% \begin{minted}{python}
% FHN = DynamicSystem2D(
%     model_FHN,
%     params=dict(delta=0.1, a=1)
% )
% \end{minted}

3. Построим бифуркационную диаграмму по параметру $a$ (Рис. \ref{fig:fhn_bifurcation_diagram}).

\begin{figure}[H]
    \centering
    \includegraphics[width=0.6\textwidth]{fhn_bifurcation_diagram.png}
    \caption{Бифуркационная диаграмма системы FHN.}
    \label{fig:fhn_bifurcation_diagram}
\end{figure}


4. Построим типичные фазовые портреты для различных значений $a$ по бифуркационной диаграмме (Рис \ref{fig:fhn_phase_portraits}).

\begin{figure}[H]
    \centering
    \includegraphics[width=1.0\textwidth]{fhn_phase_portraits.png}
    \caption{Фазовые портреты системы FHN.}
    \label{fig:fhn_phase_portraits}
\end{figure}

Для нейронной модели ФитцХью-Нагумо при значении параметра $a = 1$ наблюдается бифуркация Хопфа.
При $a < 1$ в системе существует устойчивый предельный цикл, и равновесие $(\overline{x}, \overline{y})$ неустойчиво.
При $a > 1$ предельный цикл исчезает, и равновесие $(\overline{x}, \overline{y})$ становится устойчивым.



\newpage
\subsection{Климатическая модель Зальцман-Николис (SN)}

$$
\begin{cases}
    \dot{x} = y - x \\
    \dot{y} = -a x + b y - x^2 y
\end{cases}
$$

1. Найдем равновесия системы.

\begin{multline*}
    \begin{cases}
        y - x = 0 \\
        -a x + b y - x^2 y = 0
    \end{cases}
    \implies
    \begin{cases}
        x = y \\
        x(x^2 + a - b) = 0
    \end{cases}
    \implies \\
    \implies
    \begin{cases}
        (\overline{x}_1, \overline{y}_1) = (0, 0) \\
        (\overline{x}_2, \overline{y}_2) = (\sqrt{b - a}, \sqrt{b - a}) \\
        (\overline{x}_3, \overline{y}_3) = (-\sqrt{b - a}, -\sqrt{b - a})
    \end{cases}
\end{multline*}

2. Исследуем равновесия на устойчивость.

$$
J(x, y) = \begin{pmatrix}
    -1 & 1 \\
    -a - 2xy & b - x^2
\end{pmatrix}
$$

$$
|J(x, y) - \lambda E| = 0
$$

$$
\lambda_{1,2}(x, y) = \frac{\pm \sqrt{-4a + b^2 - 2b x^2 + 2b + x^4 - 2x^2 - 8xy + 1} + b - x^2 - 1}{2}
$$

Далее с помощью численной процедуры исследуем поведение собственных чисел якобиана в равновесных точках при различных значениях параметра $a$.

% \begin{minted}{python}
% SN = DynamicSystem2D(
%     model_SN,
%     params=dict(a=1, b=2)
% )
% \end{minted}

3. Построим бифуркационную диаграмму по параметру $a$ (Рис. \ref{fig:sn_bifurcation_diagram}). Изобразим предельные циклы и равновесия:

$$
\begin{cases}
    (\overline{x}_1, \overline{y}_1) = (0, 0) \\
    (\overline{x}_2, \overline{y}_2) = (\sqrt{b - a}, \sqrt{b - a}) \\
    (\overline{x}_3, \overline{y}_3) = (-\sqrt{b - a}, -\sqrt{b - a})
\end{cases}
$$

\begin{figure}[H]
    \centering
    \includegraphics[width=0.6\textwidth]{sn_bifurcation_diagram_single.png}
    \caption{Бифуркационная диаграмма системы SN.}
    \label{fig:sn_bifurcation_diagram}
\end{figure}


4. Построим типичные фазовые портреты для различных значений $a$ по бифуркационной диаграмме (Рис. \ref{fig:sn_phase_portraits}).

\begin{figure}[H]
    \centering
    \includegraphics[width=1.0\textwidth]{sn_phase_portraits.png}
    \caption{Фазовые портреты системы SN.}
    \label{fig:sn_phase_portraits}
\end{figure}

Для климатической модели Зальцман-Николис при значении параметра $a = 1$ наблюдается бифуркация Хопфа: равновесия $(\overline{x}_2, \overline{y}_2)$, $(\overline{x}_3, \overline{y}_3)$ становятся неустойчивыми.
При $a < 1$ в системе существуют неустойчивые предельные циклы вокруг устойчивых фокусов $(\overline{x}_2, \overline{y}_2)$, $(\overline{x}_3, \overline{y}_3)$.
При $a > 1$ фокусы становятся неустойчивыми, и предельные циклы фокруг них исчезают.
При $a = 2$ наблюдается бифуркация типа ``вилка'': происходит слияние седла $(\overline{x}_1, \overline{y}_1)$ с неустойчивыми фокусами $(\overline{x}_2, \overline{y}_2)$, $(\overline{x}_3, \overline{y}_3)$.
При всех рассмотренных значениях параметра $a$ в системе сохраняется внешний устойчивый предельный цикл.

\input{lab3/lab3_part2_include}
\graphicspath{ {./lab4/images/} }

\newpage
\section{Лабораторная работа № 4: Исследование стохастического возбуждения колебаний в ФХН-модели}
\subsection{Сравнение стохастических решений}

Проведем сравнение стохастических решений для $a = 1.1$, $a = 1.01$

(а) Построим графики $x(\varepsilon)$, пропустив переходный процесс $T = 100$ (Рис. \ref{fig:excitement_amplitude}).

\begin{figure}[htbp]
    \centering
    \includegraphics[width=1\textwidth]{excitement_amplitude.png}
    \caption{Амплитуда координаты $x$ при различных значениях $\varepsilon$.}
    \label{fig:excitement_amplitude}
\end{figure}

Локализуем зону стохастического возбуждения на основе 10 симуляций (Рис. \ref{fig:excitement_zone}).

\begin{figure}[htbp]
    \centering
    \includegraphics[width=1\textwidth]{excitement_zone.png}
    \caption{$\varepsilon$-зона стохастического возбуждения.}
    \label{fig:excitement_zone}
\end{figure}


(б) Построим временные ряды при различных значениях $\varepsilon$: с возбуждением и без возбуждения (Рис. \ref{fig:time_series_1_1}, \ref{fig:time_series_1_01}).

\begin{figure}[htbp]
    \centering
    \includegraphics[width=1\textwidth]{time_series_a=1,1.png}
    \caption{Временные ряды для $a=1.1$.}
    \label{fig:time_series_1_1}
\end{figure}

\begin{figure}[htbp]
    \centering
    \includegraphics[width=1\textwidth]{time_series_a=1,01.png}
    \caption{Временные ряды для $a=1.01$.}
    \label{fig:time_series_1_01}
\end{figure}



\newpage
\subsection{Доверительные эллипсы в допороговой и надпороговой зонах}

Построим доверительные эллипсы в допороговой и надпороговой зонах для $a = 1.1$, $a = 1.01$ (Рис. \ref{fig:confidence_ellipse_1_1}, \ref{fig:confidence_ellipse_1_01}).

В режиме без возбуждения доверительные эллипсы полностью ложатся в допороговую зону, и траектории совершают лишь небольшие колебания вокруг устойчивого равновесия.
В режиме возбуждения эллипсы захватывают надпороговую зону, и шум заставляет траектории выходить из окрестности равновесия на предельный цикл.
% Это особенно хорошо видно при $a = 1.01$.

\begin{figure}[htbp]
    \centering
    \includegraphics[width=1\textwidth]{confidence_ellipse_a=1,1_deterministic.png}
    \caption{Доверительные эллипсы для $a=1.1$.}
    \label{fig:confidence_ellipse_1_1}
\end{figure}

\begin{figure}[htbp]
    \centering
    \includegraphics[width=1\textwidth]{confidence_ellipse_a=1,01_deterministic.png}
    \caption{Доверительные эллипсы для $a=1.01$.}
    \label{fig:confidence_ellipse_1_01}
\end{figure}

\graphicspath{ {./lab5/images/} }

\newpage
\section{Лабораторная работа № 5: Исследование индуцированных шумом переходов в модели SN}
\subsection{Аттракторы детерминированной модели}
\label{sec:attractors}

(A) Изобразим 2 устойчивых равновесия при значениях параметра $a = 0.5$, $a = 7$ (Рис. \ref{fig:attractors_a}).

\begin{figure}[htbp]
    \centering
    \includegraphics[width=1\textwidth]{attractors_a.png}
    \caption{Аттракторы при $a = 0.5$, $a = 7$.}
    \label{fig:attractors_a}
\end{figure}

(B) Изобразим 2 устойчивых равновесия, а также устойчивый и неустойчивый предельный цикл при значениях параметра $a = 0.715$, $a = 0.75$ (Рис. \ref{fig:attractors_b}).

\begin{figure}[htbp]
    \centering
    \includegraphics[width=1\textwidth]{attractors_b.png}
    \caption{Аттракторы при $a = 0.715$, $a = 0.75$.}
    \label{fig:attractors_b}
\end{figure}

(C)$^*$ Изобразим 2 устойчивых равновесия, а также устойчивый и 2 неустойчивых предельных цикла при значениях параметра $a = 0.8$, $a = 0.9$ (Рис. \ref{fig:attractors_c}).

\begin{figure}[htbp]
    \centering
    \includegraphics[width=1\textwidth]{attractors_c.png}
    \caption{Аттракторы при $a = 0.8$, $a = 0.9$.}
    \label{fig:attractors_c}
\end{figure}



\newpage
\subsection{Исследование ИШП между устойчивыми равновесиями для случая A}

Исследуем и визуализируем переходы между устойчивыми равновесиями системы SN для случая A из раздела \ref{sec:attractors}, при $a = 0.5$.

(a) Локализуем бассейны притяжения равновесий $M_1$, $M_2$ (Рис. \ref{fig:basins_of_attraction}).

\begin{figure}[htbp]
    \centering
    \includegraphics[width=0.6\textwidth]{basins_of_attraction.png}
    \caption{Бассейны притяжения равновесий.}
    \label{fig:basins_of_attraction}
\end{figure}

(б) Построим доверительные эллипсы для равновесий $M_1$, $M_2$ (Рис. \ref{fig:confidence_ellipse}).
С ростом $\varepsilon$ доверительные эллипсы выходят за границы бассейна притяжения своего равновесия.

\begin{figure}[htbp]
    \centering
    \includegraphics[width=0.6\textwidth]{confidence_ellipse.png}
    \caption{Доверительные эллипсы для аттракторов при $\varepsilon = 0.1$, $\varepsilon = 0.2$.}
    \label{fig:confidence_ellipse}
\end{figure}

(в) Построим временные ряды для траекторий с началом в равновесиях $M_1$, $M_2$ (Рис. \ref{fig:time_series}).
С ростом $\varepsilon$, после достижения определенного порога, можно наблюдать индуцированные шумами переходы между равновесиями системы.

\begin{figure}[htbp]
    \centering
    \includegraphics[width=1\textwidth]{time_series.png}
    \caption{Временные ряды для траекторий при $\varepsilon = 0.1$, $\varepsilon = 0.2$.}
    \label{fig:time_series}
\end{figure}



\newpage
\subsection{Исследование ИШП между равновесиями и предельным циклом для случая C}

Исследуем и визуализируем переходы между устойчивыми равновесиями системы SN для случая C из раздела \ref{sec:attractors}, при $a = 0.9$.

(а) Построим доверительные эллипсы для равновесий $M_1$, $M_2$ (Рис. \ref{fig:cycle_confidence_ellipse}).
С ростом $\varepsilon$ эллипсы выходят за границы неустойчивого предельного цикла вокруг равновесия, что делает возможным выход траектории из окрестности равновесия на предельный цикл.

\begin{figure}[htbp]
    \centering
    \includegraphics[width=0.6\textwidth]{cycle_confidence_ellipse.png}
    \caption{Доверительные эллипсы для аттракторов при $\varepsilon = 0.02$, $\varepsilon = 0.05$.}
    \label{fig:cycle_confidence_ellipse}
\end{figure}

(б) Построим временные ряды для траекторий с началом в равновесиях $M_1$, $M_2$ (Рис. \ref{fig:cycle_time_series}).
С ростом $\varepsilon$, после достижения определенного порога, можно наблюдать индуцированные шумами переходы между равновесиями и неустойчивыми предельными циклами.

\begin{figure}[htbp]
    \centering
    \includegraphics[width=1\textwidth]{cycle_time_series.png}
    \caption{Временные ряды для траекторий при $\varepsilon = 0.02$, $\varepsilon = 0.05$.}
    \label{fig:cycle_time_series}
\end{figure}


\end{document}